% =========================================================================
% Capítulo 5: Mecánica Lagrangiana
% Tema 01: Mecánica Lagrangiana
% Problema 02: Péndulo suspendido de un carro móvil
% =========================================================================
\subsubsection{Péndulo suspendido de un carro móvil}

Un carro de masa $M$ puede deslizar sin fricción a lo largo de un riel horizontal. De un punto fijo al carro cuelga un péndulo simple: una partícula de masa $m$ unida al carro mediante una varilla rígida ideal (sin masa) de longitud $\ell$, como se muestra en la Figura \ref{fig:c5_e01_p02_pendulo_carro}. Sea $x$ la posición horizontal del carro y $\theta$ el ángulo de la varilla respecto a la vertical. No hay fricción en ninguna parte del sistema.

\begin{enumerate}
    \item Identifique los grados de libertad y las coordenadas generalizadas.
    \item Construya el Lagrangiano del sistema.
    \item Obtenga las dos ecuaciones de movimiento acopladas.
    \item Verifique que $x$ es una coordenada cíclica e interprete físicamente la cantidad conservada asociada.
\end{enumerate}

\begin{figure}[!ht]
    \centering
    \begin{tikzpicture}[>=Latex, scale=1.6]
        % --- Parámetros ---
        \def\ell{2.0}      % Longitud de la varilla
        \def\Theta{30}     % Ángulo respecto a la vertical
        \def\xC{0.8}       % Posición horizontal del carro

        % --- Riel horizontal ---
        \draw[thick] (-2.5,1.4) -- (3.5,1.4);
        \fill[pattern=north east lines, pattern color=gray!60] (-2.5,1.4) rectangle (3.5,1.6);

        % --- Carro ---
        \coordinate (Ccarro) at (\xC, 1.0);
        \draw[thick, fill=gray!20] (\xC-0.6, 0.75) rectangle (\xC+0.6, 1.25);
        \node at (\xC, 1.0) {$M$};
        \filldraw (\xC, 0.75) circle (1pt) node[below] {$O$};

        % --- Referencia horizontal x ---
        \draw[->, thin] (-2.2,1.8) -- (-1.2,1.8) node[right] {$x$};

        % --- Vertical de referencia desde O ---
        \draw[dashed, gray] (\xC, 0.75) -- (\xC, 0.75-\ell-0.3);

        % --- Varilla y masa pendular ---
        \coordinate (P) at ({\xC+\ell*sin(\Theta)}, {0.75-\ell*cos(\Theta)});
        \draw[thick] (\xC, 0.75) -- (P);
        \filldraw[fill=gray!15, thick] (P) circle (3.5pt);
        \node[right] at (P) {$m$};

        % --- Angulo theta ---
        \draw[->] (\xC, {0.75-0.6}) arc (-90:{-90+\Theta}:0.6);
        \node at ({\xC}, {0.75-0.6}) [below] {};
        \node at ({\xC+0.75*sin(\Theta/2+45)}, {0.75-0.75*cos(\Theta/2+45)}) {$\theta$};

        % --- Etiqueta longitud ---
        \node[left] at ({\xC+\ell/2*sin(\Theta)-0.1}, {0.75-\ell/2*cos(\Theta)}) {$\ell$};

    \end{tikzpicture}
    \caption{Péndulo de masa $m$ suspendido de un carro de masa $M$ que desliza sin fricción sobre un riel horizontal.}
    \label{fig:c5_e01_p02_pendulo_carro}
\end{figure}

% --- SOLUCIÓN ---------------------------------------------------------------
\begin{solucion}{Péndulo suspendido de un carro móvil}
\subsubsection*{Pregunta 1: Grados de libertad}

El sistema tiene $2$ grados de libertad: la posición del carro a lo largo del riel y la orientación del péndulo. Las coordenadas generalizadas naturales son $q_1 = x$ (posición horizontal del carro) y $q_2 = \theta$ (ángulo de la varilla respecto a la vertical).

\subsubsection*{Pregunta 2: Lagrangiano}

Posición de la masa pendular respecto a un origen fijo, con $O$ (el pivote en el carro) a altura de referencia $y=0$:
$$X = x + \ell\sin\theta, \qquad Y = -\ell\cos\theta$$

Velocidades:
$$\dot{X} = \dot{x} + \ell\dot\theta\cos\theta, \qquad \dot{Y} = \ell\dot\theta\sin\theta$$

Energía cinética (carro + partícula):
$$T = \frac{1}{2}M\dot{x}^2 + \frac{1}{2}m\left(\dot{X}^2+\dot{Y}^2\right)$$

Expandiendo $\dot{X}^2+\dot{Y}^2 = \dot{x}^2 + 2\ell\dot{x}\dot\theta\cos\theta + \ell^2\dot\theta^2$:

$$T = \frac{1}{2}(M+m)\dot{x}^2 + m\ell\dot{x}\dot\theta\cos\theta + \frac{1}{2}m\ell^2\dot\theta^2$$

Energía potencial (el carro no cambia de altura):
$$V = mgY = -mg\ell\cos\theta$$

$$\boxed{L = \frac{1}{2}(M+m)\dot{x}^2 + m\ell\dot{x}\dot\theta\cos\theta + \frac{1}{2}m\ell^2\dot\theta^2 + mg\ell\cos\theta}$$

\subsubsection*{Pregunta 3: Ecuaciones de movimiento}

\textbf{Respecto a $x$:}
$$\frac{\partial L}{\partial\dot{x}} = (M+m)\dot{x} + m\ell\dot\theta\cos\theta, \qquad \frac{\partial L}{\partial x} = 0$$

$$\boxed{(M+m)\ddot{x} + m\ell\left(\ddot\theta\cos\theta - \dot\theta^2\sin\theta\right) = 0}$$

\textbf{Respecto a $\theta$:}
$$\frac{\partial L}{\partial\dot\theta} = m\ell\dot{x}\cos\theta + m\ell^2\dot\theta, \qquad \frac{\partial L}{\partial\theta} = -m\ell\dot{x}\dot\theta\sin\theta - mg\ell\sin\theta$$

Al aplicar Euler-Lagrange, los términos $\mp m\ell\dot{x}\dot\theta\sin\theta$ se cancelan, quedando:

$$\boxed{\ell\ddot\theta + \ddot{x}\cos\theta + g\sin\theta = 0}$$

Estas dos ecuaciones están acopladas: la aceleración del carro depende de $\theta$ y $\dot\theta$, y la aceleración angular depende de $\ddot{x}$.

\subsubsection*{Pregunta 4: Coordenada cíclica y cantidad conservada}

El Lagrangiano no depende explícitamente de $x$ (solo aparece $\dot{x}$), por lo que $x$ es una \textbf{coordenada cíclica}. Su momento canónico conjugado se conserva:

$$p_x = \frac{\partial L}{\partial\dot{x}} = (M+m)\dot{x} + m\ell\dot\theta\cos\theta = \text{constante}$$

\paragraph{Interpretación física} $p_x$ es exactamente el momento lineal horizontal total del sistema (carro más partícula): $p_x = M\dot x + m\dot X$. Su conservación refleja que no existe ninguna fuerza externa horizontal actuando sobre el sistema completo (el riel es sin fricción y la gravedad es vertical), consistente con el principio de conservación de momento lineal en ausencia de fuerzas externas netas en esa dirección.
\end{solucion}
